\section{Full-dimensional cluster patches}
\label{sec:full-dimensional-cluster-patches}

We now exclude patches in which diffuse clusters of second-side labels have
uniformly large common-neighbor mass on the first side. Under the initial
discrepancy bound, the next proposition turns an available family of such
patches into a cube. Their eventual absence will be used twice: to exclude
the linear-budget jump in Section~\ref{sec:jump-to-large-bias-only-at-linear-budget}
and to control the residual cluster scale during patch extraction in
Section~\ref{sec:patch-tiling-under-9}.

\begin{proposition}[Full-dimensional cluster exclusion]\label{prop:full-dimensional-cluster-exclusion}
Assume \eqref{eq:source-2} on the retained sets in both orientations. Fix constants \(\zeta>0\) and \(0<\delta<\min(\eta_0,\zeta,1)/2000\). Suppose there are available patches of \textbf{one} color \(G\) as follows:
\begin{itemize}
\item a law \(\mu_i\) on the first side and a finite probability mixture \(\lambda_i\) of cluster laws \(D\) on the second side, with \(\nu_i=\mathbb E_{\lambda_i}D\); both \(\mu_i,\nu_i\) have width \(\le n^\delta\), and each \(D\) has maximum atom \(\le \exp(-n^\zeta)\);
\item within the support of each \(D\) all pairs \(y,y'\) (repetitions allowed) have \(G\)-codegree (common-neighbor mass) under \(\mu_i\) at least \(1/4+a\), \(a=n^{-\delta}\).
\end{itemize}

Then there is a cube.
\end{proposition}

\begin{proof}\leavevmode\par
\paragraph{The finite patch menu.}
In proving this we can limit patch tags \(i\) to a finite menu still available with fixed discard fraction (retain, for each allowed discard set, a witness off it). Use \(X,Y\) for first and second sides. Take
\[
 m=\lfloor n^{200\delta}\rfloor,\qquad k=\lceil n^{300\delta}\rceil,\qquad T=\lceil n^{141\delta}\rceil .
\]
We split off \(m\) bits, writing \(z\) for their exact word. Tags \(i_z\) will be independently sampled from product profiles to be fixed later.

\subsection{Projections and lists}
\label{sec:projections-and-lists}

Write \(d=n-m\) as a sum of distinct powers of two and split the residual coordinates into chunks of those sizes \(H'\). Index the bits in one chunk by the elementary binary group of order \(H'\). For a word \(s\), its syndrome is
\(u(s)=\sum_l l\,s_l\).
Define its projection to be \(s\oplus e_{u(s)}\). The projected syndrome is \(u(s)+u(s)=0\). Conversely, a zero-syndrome word has exactly the \(H'\) preimages obtained by flipping one of its bits. This uses the binary syndrome correction rule of Hamming's parity-check construction~\cite[Sections~3--4]{hamming1950}. The index 0 is also a bit position, so exactly one bit is always flipped.

Use the product projection \(\psi\) on residual words of both parities. It toggles parity by the number of chunks modulo two. In a fixed slice \(z\), call the actual odd roles with the same projected state a group. The projected even states are the sites where centers will be selected. Across a residual edge only one projected chunk changes, in at most three coordinates; a special-coordinate flip leaves the projected residual word fixed. Thus the even projected sites incident to one odd group have pairwise residual distance at most 6 in its own slice, and there is one incident site in each adjacent slice.

A changed chunk admits at most \(O(H'^3)\) neighboring projected states. Hence each even projected site belongs to at most \(O(\sum H'^3)=O(n^3)\) group neighborhoods. The size of a projection fiber is \(\prod H'\le n^{O(\log n)}\), since there are \(O(\log n)\) chunks. In particular the group sizes have this bound.

At a single actual even role \(v\), partition its ordinary neighbors by odd group. Different chunks of the flip give different groups (the changed projected chunk has opposite parity to that at \(v\)). Within one chunk the neighbor projections are \(s\oplus e_l\oplus e_{u(s)+l}\). If the syndrome \(u(s)\ne0\), multiplicities are 2; if zero, the single multiplicity is \(H'\). Thus these ordinary incidences have multiplicities \(j\ge2\), with at most one singleton (possible unit chunk). Special incidences are all singleton groups.

Use Lemma~\ref{lem:height-selection} independently in each slice, with step bound 6 among even projected sites and parameters
\[
 r=\lfloor n^{1-10\delta}\rfloor,\qquad
 \lambda=n^{10},\qquad b_0=50\delta,\qquad b=140\delta.
\]
For the height lemma's parameters \(\zeta,\sigma,1-\theta,a\), unrelated to the cluster parameters here, take \(8\delta,\sigma_g,\delta,20\delta\), with \(0<\sigma_g\ll\delta\). The sublinear-radius condition is satisfied, and the height horizon \(H\) obeys \(r+H=O(n^{1-7\delta})\). An ID \(c\) consists of a slice, a residual location and a level. Independently pre-generate \(W_c\sim\mu_{i_z}^{\otimes k}\) at every possible ID in slice \(z\), including IDs that will be absent from the position field.

At each odd group in slice \(z\), choose a mask of the cluster prior. The finite menu of permitted masks is defined as follows. Choose a subset of second labels, retain clusters assigning it mass at least \(1/2\), provided their total prior probability is at least \(1/2\), and condition each retained cluster on that subset. Normalize the mixture weights as well. The resulting aggregate is at most \(4\nu_{i_z}\); codegrees on the retained supports are unchanged, and cluster atoms increase by at most a factor 2. The trivial mask is permitted. The distribution over masks will be specified below. It depends only on the tags at \(z\) and its adjacent slices; conditional on tags, masks are independent across groups and independent of the positions and tuple arrays.

The randomization order is now fixed. After the tags, generate prospective positions and the independent tuple arrays, then the group masks. The eligibility rule below uses these data but no activations. Afterwards generate independent activations and site-level selection ties. Only after the centers have been selected do we draw clusters and then labels. Uniform ties are pre-randomized separately at every site-level.

On successful heights, own-slice IDs in a group neighborhood number at most \(T\). Indeed the incident projected sites have heights differing by at most one, and a crowd ball about one representative at each occupied height contains all choices at that height. There is one further ID per special neighboring slice. Every group therefore sees at most \(T+m\) distinct IDs.

\paragraph{Tests for hypothetical lists.}
Consider at a group a hypothetical list \(\mathcal D\) of distinct involved IDs, at most \(T\) of its own slice and exactly one from each adjacent slice. Write \(F\) for the set of labels hitting (in \(G\)) all tuple entries there; use \(F_{-c}\) on deleting one of the tuples. For the masked cluster mixture put
\[
 A(F)=\mathbb E_D D(F)^2 .
\]
Test the following bounds:
\begin{equation}
 A(F)\ge e^{-2k|\mathcal D|},\quad
 A(F)/A(F_{-c})\ge
 \begin{cases}e^{(-\log4+.4a)k}& c\ \text{own},\\e^{-2k}&c\ \text{external}.\end{cases}
 \label{eq:source-4}
\end{equation}
We claim that, conditional on positions, masks and tags, a fixed list fails these tests with probability at most \(\exp(-c_5a^2k)\), for an absolute \(c_5>0\). Only its independent tuple entries remain random. Fix one order of its tuples, and for each tuple also use the order obtained by putting it last. Thus at most \(|\mathcal D|+1\) orders suffice for the absolute bound and all deletion bounds.

After a set \(J\) of second labels has survived the preceding hits, exposing the next first label \(x\) multiplies \(A(J)\) by
\[
 q(x)=\mathbb E_{D\ {\rm tilted\ by}\ D(J)^2}
              d_G(x;D|_J)^2.
\]
The aggregate law of \(D|_J\) under this tilt is
\[
 \overline D_J=\frac{\mathbb E_D[D(J)1_JD]}{A(J)}
      \le\frac{4\nu_{i_z}}{A(J)}.
\]
Jensen's inequality gives \(q(x)\ge d_G(x;\overline D_J)^2\). As long as preceding ratios are at least \(.24\), the width of \(\overline D_J\) is at most
\[
 n^\delta+O(1)+2k(T+m),\qquad
 k(T+m)=O(n^{500\delta})=o(n^{\eta_0}).
\]
The initial discrepancy bound \eqref{eq:source-2} therefore implies
\(q(x)\ge1/4-O(n^{-\eta_0})\),
except with probability at most \(e^{-n^{\eta_0}/2}\). Indeed, conditioning an incoming first law on a larger exceptional set would still fit that discrepancy bound and contradict the corresponding degree estimate.

For an own-slice entry, the codegree hypothesis additionally gives, at every such prefix,
\[
 \mathbb E_{x\sim\mu_{i_z}}q(x)
 =\mathbb E_{D\ {\rm tilted\ by}\ D(J)^2}
   \mathbb E_{y,y'\sim D|_J}
       \mu_{i_z}(N_G(y)\cap N_G(y'))
 \ge1/4+a.
\]
Clip the logarithmic increment below at \(\log(.24)\). Its typical lower bound is \(-\log4-O(n^{-\eta_0})\). Its expected surplus is at least \(.7a\) for large \(n\): use the last mean bound and \(\log v-\log u\ge v-u\) for \(0<u\le v\le1\), with the rare exceptional values charged separately. Thus the conditional mean increment for an own tuple is at least \(-\log4+.7a\).

Expose the \(k\) independent entries of the tuple placed last. The clipped log increments are bounded, so their deviation below \(( -\log4+.4a)k\) has probability \(\exp(-\Omega(a^2k))\). Stop this calculation if a preceding regularity bound fails and continue with fixed good increments after the stop; the stopped-prefix failures are bounded by the discrepancy exceptions already estimated. For an external last tuple the typical lower bound gives the weaker threshold \(-2k\); processing a whole list gives the absolute threshold \(-2k|\mathcal D|\). A union over the stated orders proves the fixed-list claim.

\paragraph{Eligibility and local validity.}
For each group, enumerate the hypothetical lists of the prescribed form using present IDs in its incident even-site balls, at any levels. A failed list is one that violates \eqref{eq:source-4}. Choose a maximal pairwise disjoint family of failed lists and mark every ID in that family as forbidden at each incident even site where it could be selected. The enumeration and maximal-family rule can be fixed deterministically. Every failed list meets the marked family by maximality; thus a list consisting entirely of eligible IDs passes the tests.

With probability \(1-o(1)\), all radius-\(r\) position counts at all site-levels are within relative error \(.002\) of \(\lambda\). On these position bounds, the number of lists per group is at most
\(L_n=\exp(O((m+T)\log n))\).
After fixing masks and tags, failures on disjoint lists use disjoint tuple entries and are independent. If \(p_n=e^{-c_5a^2k}\), the chance of \(n\) disjoint failures at one group is at most
\[
 L_n^n p_n^n=\exp(-\Omega(na^2k)),
 \qquad (m+T)\log n=o(a^2k).
\]
This bound also permits the union over all groups. A site lies in \(O(n^3)\) group neighborhoods. Fewer than \(n\) failed lists, each of size at most \(T+m\), at every such group remove at most
\(O(n^4(T+m))=o(\lambda)\)
IDs from its ball. Hence at least \(.99\lambda\) choices remain at every site-level, with high probability.

Now activate prospective centers independently and apply the height rule. Its largest-scale estimate gives simultaneous success in all slices on these eligibility sizes. Eligibility may depend on other slices; the height bound is uniform over all fixed eligible sets of the required sizes. The actual lists then consist of eligible IDs and pass \eqref{eq:source-4}.

Let \(\mathcal V_g\) be the following local event for a group \(g\): its incident position counts obey the stated bounds at every level; all eligibility sets consulted by the needed height rules have size at least \(.99\lambda\); the long-truncated rules select legitimate active IDs at levels less than \(H\), at sites and heights that are not bad; the own-slice fan is at most \(T\); and its realized list passes \eqref{eq:source-4}. We have proved that all \(\mathcal V_g\)'s hold simultaneously with probability \(1-o(1)\). On a local failure the later kernels use zero subprobabilities or the specified local fallback.

We record an input domain for these events and their later recomputations. Relative to the projected state of \(g\), its queried even sites are within residual distance 3 and special-coordinate distance 1. A height calculation remains in its queried slice and consults sites within residual distance \(12H\). Eligibility at one of those sites reads only the failed-list tests at incident groups, whose projected states are at distance at most 3 and whose slices differ by at most one special coordinate. Such a test reads tuples and positions in the radius-\(r\) balls of that group's incident sites, and its mask. A mask uses tags one additional special step away. Thus a group calculation reads only tags, tuple and position fields, masks, activations and ties within special-coordinate distance 3 and residual distance \(r+12H+9\).

An even star combines groups at projected distance at most 3 and special-coordinate distance at most 1 from it. For all group and star calculations below, including every candidate-tuple replacement, we may therefore use the common domain
\[
 d_H(z,z_0)\le4,\qquad d_H(t,t_0)\le R_{\rm loc},\qquad
 R_{\rm loc}=4r+40H+40\lceil\log_2 n\rceil.
\]
All levels are included. The final term covers passage between actual and projected words. Replacing a tuple only reevaluates the same list, marking, height and kernel functions on this fixed primitive domain; it does not initiate a recursive search through other eligibility rules. In particular heights remain slice-local. This domain will justify the independence of sufficiently separated calculations.

\subsection{Squared tilt, masks and presentations}
\label{sec:squared-tilt-masks-and-presentations}

On validity, sample at each group independently (given all preceding data) a masked cluster \(D\) tilted by \(D(F)^2\), further restricted to
\[
 D(F)\ge e^{-1.5 k|\mathcal D|},\qquad
 D(F)/D(F_{-c})\ \ge\
 \begin{cases}e^{(-\log2+.08a)k}&c\ \text{own},\\ e^{-1.2k}&c\ \text{external}.\end{cases}
\]
Let \(h_F\) be the probability retained by these restrictions under the squared tilt. For any ratio threshold \(t_c\),
\[
 \Pr_{D(F)^2\text{ tilt}}\!\left\{\frac{D(F)}{D(F_{-c})}<t_c\right\}
 \le t_c^2\frac{A(F_{-c})}{A(F)}.
\]
Substitution from \eqref{eq:source-4} bounds this by \(e^{-.24ak}\) for an own tuple and \(e^{-.4k}\) for an external tuple. The absolute-mass restriction fails with probability at most \(e^{-k|\mathcal D|}\). Their union is \(o(1)\), so \(h_F\ge1/2\).

Conditional on the chosen cluster, every odd role in the group has row law \(D|_F\). The reference labels are independent given the clusters. For a prescribed list \(\mathcal D\), let \(p_{\mathcal D}\) be the marginal law obtained by averaging this row over the restricted cluster draw, and set it to zero when \eqref{eq:source-4} fails. The restriction costs at most 2; cancelling one factor \(D(F)\) in the squared tilt and using its aggregate bound gives
\[
 p_{\mathcal D}\le8e^{2k(T+m)}\nu_{i_z}.
\]

Call all data through the actual center selections the pre-cluster history. It includes tags, positions, tuple arrays, masks, activations and ties. For an odd role \(b\), write \(p_b^W\) for the cluster-averaged row at its realized list, multiplied by \(1_{\mathcal V_g}\) for its group. The superscript does not condition only on \(W\); it denotes this full pre-cluster history. The same cap holds for \(p_b^W\).

We also define a reference sampling law at every pre-cluster history, without requiring global success. Draw clusters independently at the locally valid groups by the stated rule, then draw role labels independently conditional on those clusters. Invalid groups use independent local defaults. Later sublikelihoods retain \(\mathcal V_g\) for every group they query.

We now specify masks and bound tag-conditional means. For fixed neighboring tags and a deterministic configuration (the \(\mathcal D\) inputs with iid respective tuples), the mean of \(p_{\mathcal D}\) depends only on the number \(t\le T\) of own IDs in addition to the tag/mask laws. Choose a distribution of masks so these hypothetical means, for \textbf{all} \(t\), are each bounded by \(O(T+1)\nu_{i_z}\) pointwise. Indeed for nonnegative prices \(c_t(y)\) put \(S=\sum_{t,y}c_t(y)\nu_{i_z}(y)\) and discard labels with \(\sum_t c_t(y)>10S\). One can mask as specified to retain only cheap labels (they have \(\nu_{i_z}\)-mass at least \(.9\), also when \(S=0\), so at most \(.2\) of the prior on clusters has cheap mass below half); summed cost is \(\le 10(T+1)S\). Separation in the finite menu proves feasibility. Fix such strategies as lookups, independent of the choices of tag profiles.

\paragraph{From a prescribed list to the actual selection.}
Let \(\widehat p_b=\mathbb E[p_b^W\mid(i_z)_z]\), averaging the pre-cluster randomness at fixed tags. We claim
\(N\max\widehat p_b\le e^{.1m}\).
Condition first on the own-slice positions of the group of \(b\). On its position gate there are at most \(n^{O(T)}\) possible own lists: enumerate all subsets of at most \(T\) names in its incident balls. This enumeration depends only on positions. In each adjacent slice enumerate one external ID, requiring it to be present.

Fix the tuple and mask arrays. For an external queried slice \(z'\), a specified present ID \(c\) at level \(\ell\), and its position field \(P_{z'}\), let \(w_{\ell,c}(P_{z'})\) be the supremum, over fixed legal eligibility sets, of the conditional activation-and-tie probability that this ID is selected on the queried slice's size and height gates. In this upper bound discard all other components of the group event \(\mathcal V_g\). This supremum removes dependence on the particular tuple and mask values used to form eligibility. At level 0 uniform eligible-active selection gives the pointwise bound
\(w_{0,c}\le1/(.99\lambda)\).
At a higher level selection requires positive height, so the forced-center estimate in Lemma~\ref{lem:height-selection} gives
\[
 \mathbb E[w_{\ell,c}(P_{z'})\mid c\text{ present}]
       \le e^{-n^{c_0}}
\]
for some fixed \(c_0>0\).

Conditional on all positions, the activation and tie fields for the external queries are independent because they lie in distinct slices. Their bounds \(w_{\ell,c}\) depend only on the respective position fields. Those fields are independent, also with one center forced present in each slice. Hence the expectations of the bounds multiply, although the original eligibility rule uses cross-slice data. Summing named choices in one external slice, with their presence probabilities, gives
\[
 \sum_{\ell,c}\Pr(c\text{ present})
       \mathbb E[w_{\ell,c}\mid c\text{ present}]
 \le 1/.99+H\lambda e^{-n^{c_0}}=1/.99+o(1).
\]
We have ignored any requirement to present the enumerated own list, which only enlarges the probability.

The resulting weights are independent of tuple and mask values. For each deterministic list, we may now average \(p_{\mathcal D}\) using its hypothetical mean bound. Summing the own lists and external vectors gives
\[
 \widehat p_b\le
 n^{O(T)}(1/.99+o(1))^m O(T+1)\nu_{i_z}.
\]
Its normalized logarithm is
\[
 O(T\log n)+m\log(1/.99+o(1))+n^\delta+O(\log(T+1))<.1m.
\]
This proves the claim and explains why eligibility had to be fixed before activation.

\subsection{Internal weights and load transfer}
\label{sec:internal-weights-and-load-transfer}

Fix an actual even role \(v\) in slice \(z\) and a specified present candidate ID \(c\) in its center ball. Hold fixed the tags, positions, masks, activations and choice randomizers, and all tuples other than \(W_c\). The candidate tuple \(w=W_c\) has prior \(\mu_{i_z}^{\otimes k}\).

Define \(L_w\) to be the sublikelihood of the ordered odd-neighbor labels: require the position-count gate at \(v\), the selection of \(c\) there, and every neighboring group event \(\mathcal V_g\), all recomputed with \(w\). On this gate use the reference law that draws one newly tilted cluster in each group and then independent queried labels within that cluster. Thus \(L_w\) integrates the clusters; it does not fix any realized cluster or impose later injectivity constraints. We claim that
\begin{equation}
 L_w\ \le\ \exp((\log 2-.06 a)kn)\,Q
 \label{eq:source-5}
\end{equation}
for a reference probability \(Q\) independent of \(w\).

To prove this comparison, first prescribe the ID list at one neighboring group. Delete \(c\) from that list and keep the other tuples and the mask fixed. Reference its \(j\) queried labels by the masked cluster prior tilted by \(D(F_{-c})^2\), without the extra restrictions, followed by independent samples from \(D|_{F_{-c}}\).

Before integrating the cluster, the likelihood ratio on actual support is
\[
 h_F^{-1}\frac{A(F_{-c})}{A(F)}
       \left(\frac{D(F)}{D(F_{-c})}\right)^{2-j}
 \le2\frac{A(F_{-c})}{A(F)}
       \left(\frac{D(F)}{D(F_{-c})}\right)^{2-j}.
\]
The three factors come from the retained tilt mass, the two tilt normalizers, and the \(j\) conditional label densities. For an ordinary group with \(j\ge2\), the deletion bound in \eqref{eq:source-4} and the lower ratio gate give
\[
 2\exp\bigl((j\log2-.08aj-.24a)k\bigr)
       \le2e^{(\log2-.08a)kj}.
\]
For a singleton use \(D(F)/D(F_{-c})\le1\), giving at most \(2e^{2k}\). There are at most \(m+1\) singleton groups.

The realized lists may change with \(w\). At each group average the deletion references uniformly over all possible name lists containing \(c\). This probability law is independent of \(w\), and dominates any one such reference at a cost at most \(\exp(O((m+T)\log n))\) on the position gate. Nonexistent or undefined references use fixed independent fallbacks. Taking the product over groups yields a reference \(Q\) for which
\[
 \log(L_w/Q)\le(\log2-.08a)kn
       +O(km)+O(n(m+T)\log n)+O(n)
       \le(\log2-.06a)kn.
\]
The last inequality uses \(m=o(an)\) and \((m+T)\log n=o(ak)\), proving \eqref{eq:source-5}.

On the true local gates with \(c\) actually selected, test
\[
 Z=\int L_w(\boldsymbol y)\,d\mu_{i_z}^{\otimes k}(w)\ \ge e^{-.01 a k n}Q(\boldsymbol y),\qquad Z>0
\]
at the observed star data \(\boldsymbol y\). Under reference odd sampling the unconditional probability of such a true-gated failure is \(\le e^{-.01 a k n}\) per candidate test with the conditioning just described, by predictive mass. On passing, use the resulting posterior proportional to \(L_w\,d\mu_{i_z}^{\otimes k}\). It is supported on tuples of common neighbors, and has normalized joint atom cap (relative to uniform \(k\)-fold sampling on original \(X\))
\[
 \exp((\log 2-.04a)kn)
\]
using \(n^\delta=o(an)\). Let \(\overline p\) be the average coordinate marginal of this posterior and put
\[
 B=\{x:N\overline p(x)>e^{(\log2-.02a)n}\}.
\]
Then \(|B|/N\le e^{-(\log2-.02a)n}\). If \(q=\lceil(1-.01a)k\rceil\), the joint tuple cap and a union over coordinate subsets give
\[
 \Pr\{\#\{j:W_j\in B\}\ge q\}
 \le2^k e^{(\log2-.04a)kn}(|B|/N)^q
 \le e^{-\Omega(akn)}.
\]
Indeed the coefficient gained in the exponent is at least
\(.02-.01\log2+.0002a>0\), while \(k\log2=o(akn)\).
Thus the expected light-coordinate fraction is at least \(.01a(1-o(1))\). For an absolute \(c_6>0\), the light part of \(\overline p\) has mass at least \(c_6a\).

Normalize this light part to obtain the even row \(p_v^X\). In estimates set it to zero off the true local gates or predictive success. It is supported on common neighbors, and
\[
 N\max p_v^X\le e^{(\log2-.01a)n},
\]
since normalizing costs \(1/(c_6a)\) and \(\log(1/a)=o(an)\).

\paragraph{Conditional means and tag balance.}
Let \(\widehat p_v^X\) be the mean of this row over the raw local randomizations and reference neighbor draws, with tags fixed. Its normalized cap is at most \(e^{.1m}\). For a specified candidate \(c\), the subdensity \(L_w\) includes exactly the event that \(c\) is selected and the neighboring groups are locally valid. Its data marginal is \(Z\). Consequently, as measures in the candidate tuple,
\[
 \int Z(\mathbf y)
       \frac{\mu_{i_z}^{\otimes k}(dw)L_w(\mathbf y)}{Z(\mathbf y)}
          \,d\mathbf y
 =\mu_{i_z}^{\otimes k}(dw)\int L_w(\mathbf y)\,d\mathbf y
 \le\mu_{i_z}^{\otimes k}(dw),
\]
with zero terms at \(Z=0\). The last integral is a gate probability. Coordinate averaging and light-part normalization cost at most \(1/(c_6a)\). On the position gate there are only polynomially many possible candidates at \(v\). Summing their bounds and using \(N\max\mu_i\le e^{n^\delta}\) proves the claimed cap.

Both \(\widehat p_v^X\) and \(\widehat p_b\) depend only on the tags in the declared bounded special-coordinate neighborhood. Each row is supported on labels in its own patch.

Choose independent tag profiles, one per slice, so that the expected odd and even mean rows, averaged over their respective roles within that slice, are bounded coordinatewise by \(O(1)/N\). For fixed profiles elsewhere, price separation makes these bounds feasible for one slice: availability supplies a tag whose two supports avoid expensive labels, regardless of the remaining draws. The mask lookups and all other rules conditional on realized tags have already been fixed, independently of profile probabilities. The response sets are therefore nonempty, convex and have closed graph on the finite profile simplices. Lemma~\ref{lem:balanced-profiles} gives simultaneous profiles.

Under these product profiles, apply Lemma~\ref{lem:scattered-moments} to the normalized tag-conditional means. Their cap is \(e^{.1m}\). Two roles are near for this calculation when their special words lie within a sufficiently large fixed distance, so the near fraction is at most \(n^{O(1)}2^{-m}\). At separated words their mean functions use disjoint tags. The product comparison is thus independence, and
\(n\,n^{O(1)}2^{-m}e^{.1m}=o(1)\).
With high probability their averaged normalized means are bounded simultaneously for all labels. Fix such typical tags.

Next expose the pre-cluster history at these tags. Declare actual parity roles residual-near when their projected words are within \(2R_{\rm loc}\), enlarging the fixed constant if necessary. The primitive domains for separated roles are disjoint. Since \(R_{\rm loc}=O(n^{1-7\delta})\), its ball volume is \(\exp(o(an))\); the \(m\) special bits and the \(O(\log n)\) projection changes also cost only this amount. Thus the near fraction is
\[
 f_{\rm res}\le2^{-n}\exp(o(an)).
\]
For a fixed second label, the normalized row \(Np_b^W(y)\) has cap \(8e^{2k(T+m)+n^\delta}\), which is \(e^{o(n)}\). Its mean is \(N\widehat p_b(y)\). At separated roles the raw local experiments are independent given tags, even when events elsewhere fail, because all rules were extended locally. Scattered moments and the bounded average of these comparison means give odd column sums \(o(1)\) with high probability. Intersect this event with the already proved simultaneous pre-cluster validity.

On these good prehistories, draw the independent group clusters. The conditional individual laws \(D|_F\) then have column sums \(\le\theta_0\) with high probability, where \(\theta_0\) is as in the clock lemma. Indeed their sums have the preceding small expectations; the maximum contribution from a group is
\[
 n^{O(\log n)}\,2\exp(-n^\zeta+1.5k(T+m)),
\]
even including its entire group size. This is \(\exp(-\Omega(n^\zeta))\) by \(500\delta<\zeta\). Bounded-summand exponential concentration for independent groups beats the union over columns: if \(L_*\) is the displayed summand bound, a column sum with conditional mean \(\bar s=o(1)\) exceeds \(\theta_0\) with probability at most \(\exp(((e-1)\bar s-\theta_0)/L_*)\), by exponential Markov (use \(e^x\le1+(e-1)x\) on \([0,1]\)). On this success use a clock-lemma injection of all odd labels with joint upper domination through \(n^2\) queries (no additional predicates needed).

\paragraph{Transfer to even rows.}
At a successful cluster history the clock law compares each queried star, and any collection of at most \(n\) separated stars, to independent odd-row sampling conditional on those realized clusters. Retain the prehistory and cluster-column success indicators while applying this comparison; it queries at most \(n^2\) odd rows. Once the comparison is made, remove the nonlocal cluster-column indicator and integrate the independent group clusters. This gives exactly the product-by-groups reference law used to define \(L_w\), with the local gates still present.

For one star, this order of integration permits the true-gated predictive bound above to be used. Summing it over all \(2^{n-1}\) even roles and their polynomially many possible candidates gives \(o(1)\), because its exponent is \(.01akn\gg n\). Thus predictive failures occur with total probability \(o(1)\) on the entering successes.

For the even column estimate, use normalized weights \(Np_v^X(x)\) with these success indicators. Their cap is \(e^{(\log2-.01a)n}\), and
\[
 nf_{\rm res}e^{(\log2-.01a)n}=o(1).
\]
At retained separated rows, first make the same clock comparison and integrate clusters in the same order. Then remove nonlocal prehistory restrictions. The true local gates remain in each likelihood, and the declared primitive domains are disjoint. Integrating those local fields therefore gives the product of the fixed-tag comparison means \(N\widehat p_v^X(x)\). Their average is bounded by the chosen tags. Lemma~\ref{lem:scattered-moments} yields even column sums \(O(2^n/N)=o(1)\), hence at most 1, with high probability on the successes.

On their joint success every even row has mass 1 and consists of common neighbors of the odd injection. Hall's theorem supplies distinct even representatives, completing the cube and the proposition.

\end{proof}

\begin{corollary}[Eventual absence of cluster witnesses]\label{cor:eventual-cluster-absence}
After stabilization, cluster witnesses from Proposition~\ref{prop:full-dimensional-cluster-exclusion} are eventually absent for every fixed rational admissible parameter choice, in both orientations and both colors.
\end{corollary}

\begin{proof}
We may stabilize (for countably many fixed choices of these parameters) and hence assume no such cluster witnesses at all on the retained sets, either orientation or color, since availability gives the just-proved contradiction. More precisely we take in particular all fixed rational choices meeting the parameter conditions, and for each such choice we have eventual absence. Further \(o(N)\) discards do not affect previous availability conclusions.
\end{proof}
